\chapter{现代人工智能方法：表示学习、关系建模与序贯决策}
\label{ch:lower-ml-deep}

怎样把一个非线性预测写成许多简单运算的组合，又怎样求出每个参数应当调整的方向？本章先完整计算一个两层网络，再推广到矩阵形式；卷积则展示同一组参数如何在多个位置重复使用。

树模型通过划分特征空间表示非线性关系，神经网络通过复合函数学习表示。两者组织近似的方式不同，适用性取决于数据结构、样本数量和所需的预测任务。

\section{计算图、张量与反向传播}

\MFQLead{先算一个两层网络}
取标量 $h=\operatorname{ReLU}(ax+b)=\max(0,ax+b)$、$\hat y=ch+d$、$L=(\hat y-y)^2/2$，参数为 $x=2,y=1,a=1,b=-1,c=2,d=0$。前向得到 $ax+b=1,h=1,\hat y=2,L=1/2$。设误差 $e=1$，因激活输入严格为正，链式法则给
\[
 \partial_cL=eh=1,\quad \partial_dL=e=1,\quad
 \partial_aL=ecx=4,\quad \partial_bL=ec=2.
\]
学习率为 $0.1$，同时用旧参数梯度更新，得到 $(a,b,c,d)=(0.6,-1.2,1.9,-0.1)$。新激活输入为零，输出 $-0.1$、损失 $0.605$，反而上升。梯度方向只保证足够小步长的局部下降；沿一步跨过 ReLU 边界尤其不能假定下降。取学习率 $0.01$ 可得到新输出 $1.781$、损失约 $0.3053$。notebook 逐项展示梯度并扫描步长，再用一个小网络学习带噪非线性函数。

下面把同样的链式法则写成矩阵形式，以同时处理多个输入与隐藏单元。

神经网络是参数化复合函数 $f_\theta$。以两层网络为例：
\[
z_1=W_1x+b_1,\quad h=\phi(z_1),\quad
\hat y=W_2h+b_2,\quad
L=\ell(\hat y,y).
\]
反向传播沿计算图反向使用链式法则，重复利用已经算出的局部导数：
\[
\frac{\partial L}{\partial W_2}
=\frac{\partial L}{\partial\hat y}h^{\mathsf T},
\qquad
\frac{\partial L}{\partial W_1}
=\left(W_2^{\mathsf T}\frac{\partial L}{\partial\hat y}
\odot\phi'(z_1)\right)x^{\mathsf T}.
\]
自动微分也按这些局部规则累积导数。因此，手算小网络的中间梯度可以帮助理解框架返回的梯度；实际训练的目标则由我们选择的损失函数决定。

张量的形状说明每一轴组织哪些观测。金融序列常写成 $X\in\mathbb R^{B\times T\times F}$，分别表示批次、时间和特征。若把资产维与时间维交换，代码仍可能运行，但模型学习的问题已经改变。例如最后一轴含收益和成交量两个特征时，交换时间轴与特征轴就改变了哪些数在一起参与运算。

\MFQLead{激活、初始化与梯度传播}

若每层都使用线性映射，多层网络仍等价于单个线性变换。ReLU $\phi(z)=\max(0,z)$ 提供分段线性表达，但负半轴梯度为零；GELU 以平滑门控常用于 Transformer。Sigmoid 和 $\tanh$ 在绝对值大时导数接近零，长链相乘容易导致梯度消失。

初始化要稳定前向方差与反向梯度。若输入分量独立、方差为 $v$，权重方差约取 $1/n_{\rm in}$ 可避免线性层方差随宽度爆炸；ReLU 常采用 He 初始化 $2/n_{\rm in}$。该近似基于对称的零均值预激活：ReLU 将二阶矩减半；激活后均值通常不为零，因此不能把二阶矩直接称为方差。

残差连接
\[
h_{l+1}=h_l+F_l(h_l)
\]
为梯度提供恒等路径。LayerNorm 对单个样本的特征维标准化，BatchNorm 使用批次统计；金融小批次、非平稳和跨资产混合会使 BatchNorm 的运行统计难以解释，不能照搬图像默认配置。

\section{优化器、正则化和训练循环}

随机梯度下降用抽取的小批次梯度近似训练集平均损失的梯度。Momentum 累积方向，Adam 对一阶和二阶矩做自适应缩放。AdamW 将权重衰减从梯度更新中解耦。优化器改变训练动力学，不改变数据的有效信息量；在只有数十个独立月份的样本上，百万行资产面板并不等于百万个独立观察。

% Generated by tools/build_figures.py; edit figures/figure-specs.json instead.
\MFQFigure{tex/figures/generated/lower-ml-training.pdf}{本例的训练损失继续下降，验证损失却开始上升，表现出过拟合；早停使用验证结果选择迭代轮数。}{lower-ml-training}


常见正则化包括权重衰减、dropout、数据增强、早停和容量限制。早停监控的是验证窗口，不是最终测试窗口。Dropout 在训练时随机屏蔽单元，推理时必须切到 \texttt{eval()}；忘记切换会使同一输入得到随机预测。

训练可以理解为重复三步：前向计算预测与损失，反向计算梯度，再按学习率更新参数。每轮之后在独立验证窗口计算损失，观察训练误差下降是否伴随未来误差下降。选择早停轮数属于模型选择，最终测试数据不参与这一步。

\section{卷积：局部连接与参数共享}

一维卷积核 $K\in\mathbb R^{k\times F\times C}$ 在时间轴滑动：
\[
H_{t,c}=\phi\!\left(b_c+
\sum_{u=0}^{k-1}\sum_{f=1}^F K_{u,f,c}X_{t-u,f}\right).
\]
因果卷积只读取 $t$ 及以前；普通“same”卷积若对称填充，会读取未来位置。扩张卷积以间隔 $d$ 取样，使较少层覆盖更长历史。$L$ 层、核宽 $k$、扩张率 $1,2,\ldots,2^{L-1}$ 的感受野为
\[
1+(k-1)(2^L-1).
\]

% Generated by tools/build_figures.py; edit figures/figure-specs.json instead.
\MFQFigure{tex/figures/generated/lower-ml-cnn.pdf}{一维因果卷积只汇总当前及过去窗口；感受野随层数扩大，而参数在时间位置间共享。}{lower-ml-cnn}


\MFQLead{卷积的手算例}

价格变化序列为 $(1,3,2,5)$，因果核为 $(1,-1)$，按“当前减前一期”定义输出，则有效位置得到 $(2,-1,3)$。反转序列后输出变为 $(-3,1,-2)$；时间均值虽然不变，卷积保留了局部顺序。比较两种次序的输出，可以看出卷积保留了哪些局部变化。

二维卷积适合图像、相关矩阵或订单簿“价格档位 $\times$ 时间”的网格表示，但网格的邻接关系必须有经济含义。把任意因子表 reshape 成图片并不会创造空间局部性。

\MFQLead{CNN 在量化研究中的三种典型输入}

第一类是多变量时间序列：收益、成交量、波动和盘口特征沿时间排列，用一维因果 CNN 提取局部形状。第二类是限价订单簿张量：价格档、买卖方向和事件时间形成局部结构。第三类是文本或事件 token 的局部 n-gram 表示；这类任务如今常由 Transformer 接管，但 CNN 仍是速度快、延迟低的基线。

所有输入都要明确采样时钟。日线卷积中“过去 20 根 bar”较清楚；事件时间卷积中 20 个事件可能跨越毫秒或数小时，需要额外时间间隔特征。对缺失交易日补零会把“未观察”混成“收益为零”，应使用 mask 或显式缺失类别。

\begin{diagnostic}
预测时点 $t$ 的结果时，以 $t$ 为中心、包含 $t$ 之后数据的滚动窗口会使用未来信息。若整个输入窗口都已在决策前观测到，窗口内部的双向运算本身并不造成泄漏。需要分别判断原始数据的可得时间与每个输出实际依赖的位置。
\end{diagnostic}

\section{分类、回归、排序与多任务学习}

方向分类使用交叉熵，收益回归常用 MSE/Huber，横截面选股更关注排序。Pairwise loss 比较同一时点资产对，Listwise loss 直接作用于整组排序。多任务学习可共享表示，同时预测收益、波动和成交概率：
\[
L=\lambda_r L_{\rm return}+\lambda_v L_{\rm vol}+
\lambda_f L_{\rm fill}.
\]
$\lambda$ 不只是调参；它定义了研究目标的相对单位。可以用验证期决策损失选权重，但不能看到测试期收益后再改。

深度集成、bootstrap、分位数损失和 conformal 方法从不同角度估计预测不确定性，各自需要相应假设。MC dropout 产生多组预测，但这些变化是否对应可靠的概率，需要通过校准或覆盖率另行考察。后续决策还要说明预测不确定性如何影响仓位。

\MFQLead{何时深度模型没有带来价值}

复杂度本身无法说明预测是否改善。可以先比较树与网络在相同时间窗口上的预测损失，再考察不同随机初始化下的变化。若讨论交易，还应使用相同的仓位与成本约定，把预测改善与更高换手带来的影响分开。

两层网络展示链式法则如何传递误差，卷积展示同一组参数如何作用于多个位置。下一章从这些运算继续研究序列依赖与 Attention。
